%               .-.
%              (o o)      SCIENCE OWL
%              | O \      (guardian of rigor)
%               \   \
%                `~~~'
%          .-------------------.
%          |  NO FABRICATION   |
%          |  + CLEAR ASSUMES  |
%          |  + CLEAN LATEX    |
%          '-------------------'
%     __        __        __        __
%  __/  \__  __/  \__  __/  \__  __/  \__
% /  RULES  \/  FLOW  \/  MATH  \/  CODE  \
% \__    __/\__    __/\__    __/\__    __/
%    \__/      \__/      \__/      \__/

%         \  |  /        \  |  /
%          \ | /          \ | /
%       ---- * ----    ---- * ----
%          / | \          / | \
%         /  |  \        /  |  \

%              ...                            
%             ;::::;                           
%           ;::::; :;                          
%         ;:::::'   :;                         
%        ;:::::;     ;.                        
%       ,:::::'       ;           OOO\         
%       ::::::;       ;          OOOOO\        
%       ;:::::;       ;         OOOOOOOO       
%      ,;::::::;     ;'         / OOOOOOO      
%    ;:::::::::`. ,,,;.        /  / DOOOOOO    
%  .';:::::::::::::::::;,     /  /     DOOOO   
% ,::::::;::::::;;;;::::;,   /  /        DOOO  
%;`::::::`'::::::;;;::::: ,#/  /          DOOO 
%:`:::::::`;::::::;;::: ;::#  /            DOOO
%::`:::::::`;:::::::: ;::::# /              DOO
%`:`:::::::`;:::::: ;::::::#/               DOO
% :::`:::::::`;; ;:::::::::##                OO
% ::::`:::::::`;::::::::;:::#                OO
% `:::::`::::::::::::;'`:;::#                O 
%  `:::::`::::::::;' /  / `:#                  
%   ::::::`:::::;'  /  /   `#              
%
%
%
%                           _,,,.._       ,_
%                        .gMMMMMMMMMp,_    `\
%                     .dMMP'       ``^YMb..dP
%                    dMMP'
%                    MMM:
%                    YMMb.
%                     YMMMb.
%                      `YMM/|Mb.  ,__
%                   _,,-~`--..-~-'_,/`--,,,____
%               `\,_,/',_.-~_..-~/' ,/---~~~"""`\
%          _,_,,,\q\q/'    \,,-~'_,/`````-,7.
%         `@v@`\\,,,,__   \,,-~~"__/` ",,/MMMMb.
%          `--''_..-~~\   \,-~~""  `\_,/ `^YMMMMMb..
%           ,|``-~~--./_,,_  _,,-~~'/_      `YMMMMMMMb.
%         ,/  `\,_,,/`\    `\,___,,,/M/'      `YMMMMMMMb
%                     ;  _,,/__...|MMM/         YMMMMMMMb
%                      .' /'      dMMM\         !MMMMMMMMb
%                   ,-'.-'""~~~--/M|M' \        !MMMMMMMMM
%                 ,/ .|...._____/MMM\   b       gMMMMMMMMM
%              ,'/'\/          dMMP/'   M.     ,MMMMMMMMMP
%             / `\;/~~~~----...MP'     ,MMb..,dMMMMMMMMM'
%            / ,_  |          _/      dMMMMMMMMMMMMMMMMB
%            \  |\,\,,,,___ _/    _,dMMMMMMMMMMMP".emmmb,
%             `.\  gY.     /      7MMMMMMMMMMP"..emmMMMMM
%                .dMMMb,-..|       `.~~"""```|dMMMMP'MMP'
%               .MMMMP^"""/ .7 ,  _  \,---~""`^YMMP'MM;
%             _dMMMP'   ,' / | |\ \\  }          PM^M^b
%          _,' _,  \_.._`./  } ; \ \``'      __,'_` _  `._
%      ,-~/'./'| 7`,,__,}`   ``   ``        // _/`| 7``-._`}
%     |_,}__{  {,/'   ``                    `\{_  {,/'   ``
%     ``  ```   ``                            ``   ``
%
%
%
%
\documentclass[11pt]{article}

% This template is designed for XeLaTeX.
\usepackage{geometry}
\geometry{
  a4paper,
  margin=25.4mm
}

\makeatletter
\def\input@path{{../Style/}}
\makeatother

\usepackage{coursework-report-core}
\usepackage{wrapfig}

% Layout debugging is provided by coursework-report-core.sty.
% Use \showlayouttrue to display page-layout frames and \showlayoutfalse for
% submission-ready output. The package option [showlayout] is also supported.
\showlayouttrue

\setlength{\parindent}{20pt}
\setlength{\parskip}{0.5em}

% Optional bibliography file for BibLaTeX/Biber.
\addbibresource{References/.bib}
\graphicspath{{Figures/}}
% Set the default level used by \smartsubsection.
% Valid values: section, subsection, subsubsection, paragraph, subparagraph
\setsmartsubsectionlevel{subsection}
% Optional depth offsets:
% \smartsubsection{Subsection}
% \smartsubsection[1]{Subsubsection}
% \smartsubsection[2]{Paragraph}

% Override these if you want a different numbering style for deeper headings.
% Example:
% \renewcommand{\paragraphnumberformat}{\thesubsection.\arabic{paragraph}}
% \renewcommand{\subparagraphnumberformat}{\theparagraph.\arabic{subparagraph}}

\fancyhead[RO]{Studentship Project - Complex Tsunami Computational Fluid Dynamics}

\begin{document}

\pagenumbering{Roman}
\tableofcontents
%\printnomenclature
\newpage
\pagenumbering{arabic}

\section*{\centering Generating Tsunami Energy-Dissipating Barriers Using Machine Learning}

\section{Core Concept}

The purpose of this studentship is to extend the \texttt{MECH0020} project by introducing greater complexity namely by using both the Shallow Water Equations (SWE) and the Navier-Stokes Equations (NSE) to develop a   

\section{Motivation}

  The 11 March 2011 T\={o}hoku tsunami will serve as the primary benchmark case study, on two accounts the first being high-quality observational data are available for validation \cite{NCEITohokuEvent,Goto2012} and second is the Sanriku/Tohoku coastline is one of the most historically tsunami-vulnerable regions in Japan, having experienced major destructive events such as the 1896 and 1933 Sanriku tsunamis, while the 2011 event itself produced extreme run-up along the same coast \cite{Fujii2011,Satake2014,Goto2012}.

\section{Action Plan}

  

\section{Validation Methods}

\section{Mathematical Model Structure}

\section{Numerical Methods}

\section{Sustainability & Optimisation}

T






































% ==============================================================================
% REUSABLE COURSEWORK SNIPPETS
% ==============================================================================

% Section-tier mapping:
% \SECTION_TIER_1{Title} -> \section{Title}
% \SECTION_TIER_2{Title} -> \subsection{Title}
% \SECTION_TIER_3{Title} -> \subsubsection{Title}
% \SECTION_TIER_4{Title} -> \paragraph{Title}
% \SECTION_TIER_5{Title} -> \subparagraph{Title}

% +----------------------------------------------------------------------------+
% | SECTION PURPOSE                                                            |
% +----------------------------------------------------------------------------+
% | Describe the scope, purpose, and expected content of this section.          |
% +----------------------------------------------------------------------------+
%
% ==============================================================================
%                                SECTION TIER 1
% ==============================================================================
% \SECTION_TIER_1{Section Title}
% \label{sec:section-label}
%
% Section content.
%
% ==============================================================================
%                            END OF SECTION TIER 1
% ==============================================================================

% Smart references:
% \smartref{sec:section-label}
% \smartref{fig:example}
% \smartref{fig:example-a,fig:example-b}
% \smartref{fig:example,tab:example}

% Configurable legacy subsection helper:
% \smartsubsection{Heading at the configured base level}
% \smartsubsection[1]{Heading one level deeper}

% Numbered equation:
% \begin{align}
%   y &= mx + c
%   \label{eq:linear-relation}
% \end{align}

% Figure:
% \begin{figure}[htbp]
%   \centering
%   \includegraphics[width=0.8\linewidth]{figures/example}
%   \caption{Example figure caption}
%   \label{fig:example}
% \end{figure}

% Subfigures:
% \begin{figure}[htbp]
%   \centering
%   \begin{subfigure}[t]{0.48\linewidth}
%     \centering
%     \includegraphics[width=\linewidth]{figures/example-a}
%     \caption{First subfigure}
%     \label{fig:example-a}
%   \end{subfigure}
%   \hfill
%   \begin{subfigure}[t]{0.48\linewidth}
%     \centering
%     \includegraphics[width=\linewidth]{figures/example-b}
%     \caption{Second subfigure}
%     \label{fig:example-b}
%   \end{subfigure}
%   \caption{Combined figure caption}
%   \label{fig:example-combined}
% \end{figure}

% Table:
% \begin{table}[htbp]
%   \caption{Example table caption}
%   \label{tab:example}
%   \begin{tabular}{lll}
%     \toprule
%     Quantity & Symbol & Unit \\
%     \midrule
%     Velocity & $v$ & \si{\metre\per\second} \\
%     \bottomrule
%   \end{tabular}
% \end{table}

% Nomenclature entry:
% \nomenclature{$v$}{Velocity, \si{\metre\per\second}}

% Algorithm:
% \begin{algorithm}
%   \caption{Example algorithm}
%   \label{alg:example}
%   \begin{algorithmic}[1]
%     \RequireWrap{Input data}
%     \EnsureWrap{Processed output}
%     \StateWrap{Long algorithm lines wrap cleanly to the next line.}
%   \end{algorithmic}
% \end{algorithm}

% MATLAB listing:
% \begin{lstlisting}[
%   style=MATLAB,
%   caption={Example MATLAB listing},
%   label={lst:matlab-example}
% ]
% x = linspace(0, 2*pi, 1000);
% y = sin(x);
% plot(x, y)
% \end{lstlisting}

% Python listing:
% \begin{lstlisting}[
%   style=Python,
%   caption={Example Python listing},
%   label={lst:python-example}
% ]
% import numpy as np
% print(np.linspace(0.0, 1.0, 5))
% \end{lstlisting}

% C++ listing:
% \begin{lstlisting}[
%   style=C++,
%   caption={Example C++ listing},
%   label={lst:cpp-example}
% ]
% #include <iostream>
%
% int main()
% {
%     std::cout << "Hello world\\n";
%     return 0;
% }
% \end{lstlisting}

% Landscape longtable:
% \begin{landscapetablepages}
%   \begin{longtable}{ll}
%     \caption{Wide landscape longtable example}
%     \label{tab:landscape-example} \\
%     \toprule
%     A & B \\
%     \midrule
%     \endfirsthead
%     \toprule
%     A & B \\
%     \midrule
%     \endhead
%     Value 1 & Value 2 \\
%     \bottomrule
%   \end{longtable}
% \end{landscapetablepages}

% Inline helpers:
% \vect{u}
% \dd
% \code{variable_name}
% \data{results/output.csv}

% Bibliography output:
% \printbibliography

\end{document}
